\section{Analisi dei principi fondamentali}

\subsection{Vincoli}

\note{Lezione 17/03/2026}
Verranno considerati sistemi finiti di punti e di corpi rigidi con vincoli che ne limitano le posizioni. Questi vincoli possono essere in moto assegnato, su cui non vi è influenza da parte del sistema.

L'insieme dei parametri che occorrono per determinare la posizione del corpo è
\begin{equation*}
    \xi_1, \xi_2, \dots, \xi_\ell.
\end{equation*}

\begin{definition}[Vincoli olonomi]
I vincoli olonomi non dipendono dalle velocità e sono espressi da disequazioni del tipo
\begin{align*}
    \varphi_\alpha(\xi_1, \xi_2, \dots, \xi_\ell, t) &= 0 \quad \alpha = 1, 2, \dots, r \note{bilateri}\\
    \psi_\beta(\xi_1, \xi_2, \dots, \xi_\ell, t) &\geq 0 \quad \beta = 1, 2, \dots, s.\note{unilaterali}
\end{align*}
\end{definition}

\begin{definition}[Gradi di libertà]
Il numero di gradi di libertà del sistema è dato da
\begin{equation*}
    n = \ell - \rho,
\end{equation*}
dove $\ell$ è il numero di parametri e $\rho$ è il rango della matrice Jacobiana dei vincoli, ovvero
\begin{equation*}
    \rho = \rank \left( \frac{\partial (\varphi_1, \varphi_2, \dots, \varphi_r)}{\partial (\xi_1, \xi_2, \dots, \xi_\ell)} \right).
\end{equation*}
Per cui si ha che il numero di gradi di libertà è dato dal numero dei parametri meno il numero di vincoli linearmente indipendenti.
\end{definition}

\subsection{Coordinate lagrangiane}
Possiamo dimostrare che esistono $n$ funzioni indipendenti $q_1, q_2, \dots, q_n$ tali che
\begin{equation*}
    \xi_i = \xi_i(q_1, q_2, \dots, q_n, t) \quad i = 1, 2, \dots, \ell.
\end{equation*}
La dipendenza dal tempo è presente solo se almeno uno di questi parametri dipende esplicitamente dal tempo.

Questo procedimento è la generalizzazione dell'espressione di $x$ e $y$ in funzione di $\theta$ per le coordinate polari.

\begin{definition}[Coordinate lagrangiane]
Questi $n$ parametri $q_1, q_2, \dots, q_n$ sono detti \emph{coordinate lagrangiane}, \emph{coordinate generalizzate} o \emph{coordinate libere}\note{libere enfatizza il fatto che siano indipendenti} del sistema.
\end{definition}

\begin{example}[Guida rettilinea ruotante]
Consideriamo un punto che può scorrere liberamente su una guida rettilinea ruotante nel piano $(O, x, y)$ che si muove con una velocità angolare $\omega$ positiva. Consideriamo la distanza $\delta=\abs{OP}$ tra il punto e l'origine. Descriviamo la posizione del punto con
\begin{equation*}
    \begin{cases}
        x = \delta \cos(\omega t) \\
        y = \delta \sin(\omega t)
    \end{cases}
\end{equation*}
avendo quindi la dipendenza dal tempo.
\end{example}

\begin{remark}
Le coordinate lagrangiane non sono univoche: esiste una corrispondenza più forte della biunivoca fra le coordinate lagrangiane e i parametri $\xi_1, \xi_2, \dots, \xi_\ell$, detta \emph{diffeomorfismo}. In particolare, se $q_1, q_2, \dots, q_n$ sono coordinate lagrangiane, allora anche $q_1', q_2', \dots, q_n'$ con
\begin{equation*}
    q_i' = q_i'(q_1, q_2, \dots, q_n, t) \quad i = 1, 2, \dots, n
\end{equation*}
sono coordinate lagrangiane.
\end{remark}

\subsection{Spostamenti infinitesimi}
Sia $\mathcal{C}_n$ uno spazio delle configurazioni reale $n$-dimensionale. Lo consideriamo come fosse uno spazio $\R^n$\note{È così solo localmente}.

Il moto del sistema è descritto fissando le funzioni
\begin{equation*}
    \begin{cases}
        q_k = q_k(t) \\
        h = 1, 2, \dots, n
    \end{cases}
    \quad t \in [t_1, t_2]
\end{equation*}
le quali hanno "incorporate" il rispetto dei vincoli bilateri. Solitamente supponiamo che queste funzioni siano di classe $C^2([t_1, t_2])$.

Introduciamo ora il concetto di \emph{spostamenti infinitesimi} o \emph{variazioni}.\note{Non ci sono approssimazioni, sotto ci sono formulazioni esatte}

Consideriamo il sistema di coordinate lagrangiane
\begin{equation*}
    P = P(q_1, q_2, \dots, q_n; t).
\end{equation*}

Possiamo calcolare la velocità del punto $P$ come
\begin{equation*}
    \vv = \dv{P}{t} = \sum_{h=1}^n \pdv{P}{q_h} \dot{q}_h + \pdv{P}{t},
\end{equation*}
in modo da definire lo \emph{spostamento infinitesimo}:
\begin{equation*}
    \dd{P} = \vv \dd{t} = \sum_{h=1}^n \pdv{P}{q_h} \dd{q}_h + \pdv{P}{t} \dd{t}.
\end{equation*}

\begin{example}[Circonferenza, versore trasversale]
\textbf{Esempio circonferenza versore trasversale etc. da completare.}
\end{example}

In un caso con una sola coordinata $s$ possiamo scrivere
\begin{equation*}
    \dd{P} = \pdv{P}{s} \dd{s} + \pdv{P}{t} \dd{t},
\end{equation*}
mentre definiamo
\begin{equation*}
    \var{P} = \pdv{P}{s} \var{s},
\end{equation*}
caso in cui non c'è la dipendenza dal tempo.

\begin{definition}[Spostamenti possibili e virtuali]
Parliamo di spostamenti \emph{reali} quando $\dd{t} \neq 0$ e di spostamenti \emph{virtuali} quando $\dd{t} = 0$.\note{Probabilmente Stevino nel '500 è stato il primo a parlare di spostamenti virtuali, con maggiore sistematizzazione da parte di Bernoulli e poi Lagrange.}
Per cui abbiamo:
\begin{align*}
    \dd{P} &= \sum_{h=1}^n \pdv{P}{q_h} \dd{q}_h + \pdv{P}{t} \dd{t} \quad &\text{spostamenti possibili} \\
    \var{P} &= \sum_{h=1}^n \pdv{P}{q_h} \delta q_h \quad &\text{spostamenti virtuali}.
\end{align*}
\end{definition}

Tornando alla particella vincolata a muoversi su una guida ruotante, se consideriamo l'istante $t + \dd{t}$, vi sarà uno spostamento possibile $\dd{P}$. Invece se consideriamo il vincolo come se fosse fisso, solitamente si dice \emph{vincolo congelato}, allora si avrà uno spostamento virtuale $\var{P}$. \textbf{inserisci figura}

Spostamenti virtuali e reali coincidono se si hanno vincoli indipendenti dal tempo.

\subsubsection{Vincoli unilateri e spostamenti ammessi}
Le coordinate lagrangiane non hanno incorporato il rispetto dei vincoli unilaterali per cui bisogna fare in modo che gli spostamenti virtuali rispettino anche questi vincoli.
\begin{equation*}
    \begin{cases}
        g_j(q_1, q_2, \dots, q_n, t) \geq 0 \\
        j = 1, 2, \dots, s
    \end{cases}
\end{equation*}

\begin{example}[Punto in una circonferenza]
Consideriamo un punto che si muove all'interno di una circonferenza di raggio $r$ con centro nell'origine.
\begin{equation*}
    \begin{cases}
        x = \rho \cos \theta \\
        y = \rho \sin \theta
    \end{cases}
\end{equation*}
Il vincolo unilaterale che abbiamo è
\begin{equation*}
    x^2 + y^2 \leq r^2 \implies \rho \leq r.
\end{equation*}

Tutti i vincoli unilaterali sono soddisfatti come disuguaglianze strette, dandoci le posizioni (o configurazioni) ordinarie
\begin{equation*}
    x^2 + y^2 < r^2 \implies \rho < r.
\end{equation*}
Invece parliamo di posizioni di confine quando almeno una delle disequazioni è verificata come uguaglianza, in questo caso avremo
\begin{equation*}
    x^2 + y^2 = r^2 \implies \rho = r.
\end{equation*}
\end{example}

Ci sono spostamenti infinitesimi consentiti, che rispettano i vincoli unilaterali, e spostamenti infinitesimi non consentiti, che violano almeno uno dei vincoli unilaterali, per esempio uscendo dalla circonferenza.

\begin{definition}[Reversibilità]
Diremo che uno spostamento virtuale è \emph{reversibile} se il suo opposto è anch'esso uno spostamento virtuale, altrimenti diremo che è \emph{irreversibile}.
\end{definition}

Tornando al disco che si espande. \textbf{inserisci figura}.

\begin{remark}
Nelle configurazioni ordinarie (tutti i vincoli unilaterali soddisfatti strettamente) gli spostamenti virtuali sufficientemente piccoli sono reversibili; nelle configurazioni di confine compaiono spostamenti virtuali irreversibili.
\end{remark}

Matematicamente:
\begin{align*}
    \text{Configurazione ordinaria:}\quad
    g_j(q_1, q_2, \dots, q_n, t) &> 0 \qquad \forall j = 1, 2, \dots, s, \\
    g_j(q_1 \pm \var{q_1}, q_2 \pm \var{q_2}, \dots, q_n \pm \var{q_n}, t) &> 0
    \qquad \forall j,\ \forall \abs{\var{q_k}} < \varepsilon.
\end{align*}

\begin{align*}
    \text{Configurazione di confine:}\quad
    \exists j^*:\ g_{j^*}(q_1, q_2, \dots, q_n, t) &= 0, \\
    g_{j^*}(q_1 + \var{q_1}, q_2 + \var{q_2}, \dots, q_n + \var{q_n}, t) &> 0
    \qquad \forall \abs{\var{q_k}} < \varepsilon, \\
    g_{j^*}(q_1 - \var{q_1}, q_2 - \var{q_2}, \dots, q_n - \var{q_n}, t) &< 0
    \qquad \text{per qualche } \abs{\var{q_k}} < \varepsilon.
\end{align*}

\subsection{Problema meccanico dei vincoli}
\begin{proposition}[Postulato delle reazioni vincolari]
Ogni vincolo può essere sostituito con un sistema di forze, dette \emph{reazioni vincolari}, senza alterare lo stato di quiete o di moto del sistema materiale.
\end{proposition}

Supporremo vincoli \emph{lisci}.
\begin{definition}[Superficie liscia]
    Diremo che una superficie è liscia, o priva di attrito, se è in grado di esercitare tutte e sole le reazioni vincolari ad essa normali e rivolte verso l'esterno, cioè verso la regione in cui il sistema può muoversi.
\end{definition}

Diremo che un sistema materiale ha vincoli privi di attrito se tutte le superfici degli ostacoli esterni e dei corpi del sistema sono lisce.

\textbf{Figura}

\begin{equation*}
    \dd{L} = \vF\cdot \dd{\vP} = \sum_{i=1}^n F_i \cdot \dd{\vP_i}.
\end{equation*}

Possiamo anche definire il lavoro virtuale
\begin{equation*}
    \delta L = \vF\cdot \var{\vP} = \sum_{i=1}^n F_i \cdot \var{\vP_i}.
\end{equation*}

Indichiamo con $\var{L}^{(V)}$ il lavoro virtuale delle reazioni vincolari.
\begin{proposition}[Vincoli lisci e lavoro virtuale]
Per un sistema olonomo con vincoli lisci si ha
\begin{equation*}
    \var{L}^{(V)} \geq 0 \quad \forall \var{P} \text{ spostamento virtuale}.
\end{equation*}
In particolare, $\var{L}^{(V)} = 0$ per gli spostamenti virtuali reversibili.
\end{proposition}

\begin{definition}
    Diremo \emph{vincoli ideali} tutti i vincoli che verificano la disuguaglianza
    \begin{equation*}
        \var{L}^{(V)} \geq 0 \quad \forall \var{P} \text{ spostamento virtuale}.
    \end{equation*}
\end{definition}

\begin{remark}
    Esistono vincoli ideali che non sono privi di attrito.
\end{remark}

Questa è l'espressione analitica del fatto che la reazione vincolare deve essere normale alla superficie di contatto.

Un sistema olonomo si dice a vincoli ideali se \note{Lezione 19/03/2026}
\begin{equation*}
    Se \var{L}^{(V)} \geq 0
\end{equation*}
per ogni spostamento virtuale effettuato a partire da una configurazione fissata.

Proprietà: Se $\mathcal{S}$ p un sistema a vincoli privi di attrito allora è a vincoli ideali.

\section{Equazioni di Lagrange}
Talvolta le equazioni che seguono sono dette di Eulero-Lagrange, ma la derivazione che faremo è propriamente di Lagrange. Risale al 1788, con la pubblicazione del \emph{Mécanique Analytique}.

Consideriamo un sistema olonomo di punti
\begin{equation*}
    \mathcal{S} = \{(\vP_1, 
    m_1), (\vP_2, m_2), \dots, (\vP_N, m_N)\}.
\end{equation*}
dove $m_i$ è la massa del punto $\vP_i$, a cui è soggetto un ssitema di forze\note{o sollecitazione}
\begin{equation*}
    \Sigma=\{(\vP_1, \vF_1), (\vP_2, \vF_2), \dots, (\vP_N, \vF_N)\},
\end{equation*}

e un sistema di coordinate lagrangiane con le relative velcità 
\begin{equation*}
    \dot{\vq} = (\dot{q}_1, \dot{q}_2, \dots, \dot{q}_n) \qquad \vq=(q_1, q_2, \dots, q_n).
\end{equation*}

La generica forza $\vF_i$ può dipendere dalle posizioni dei punti, le loro velocità, e dal tempo:
\begin{equation*}
    \vF_i = \vF_i(\vP_1, \vP_2, \dots, \vP_N; \dot{\vP}_1, \dot{\vP}_2, \dots, \dot{\vP}_N; t).
\end{equation*}

Ricordiamo che la velocità di un punto $\vP_i$ è data da
\begin{equation*}
    \vv=\dv{\vP_i}{t} = \sum_{h=1}^n \pdv{\vP_i}{q_h} \dot{q}_h + \pdv{\vP_i}{t},
\end{equation*}
per cui possiamo scrivere la forza $\vF_i$ come
\begin{equation*}
    \vF_i = \vF_i(\vP_1, \vP_2, \dots, \vP_N; \vv; t).
\end{equation*}

Calcoliamo il lavoro virtuale totale
\begin{equation*}
    \var{L} = \sum_{i=1}^N \vF_i \cdot \var{\vP}_i = \sum_{i=1}^N \vF_i \cdot \sum_{h=1}^n \pdv{\vP_i}{q_h} \var{q}_h = \sum_{h=1}^n \left( \sum_{i=1}^N \vF_i \cdot \pdv{\vP_i}{q_h} \right) \var{q}_h.
\end{equation*}
Definiamo le componenti lagrangiane della sollecitaizione
\begin{equation*}
    Q_k \defeq \sum_{i=1}^N \vF_i \cdot \pdv{\vP_i}{q_k},
\end{equation*}
ottenuiamo quindi
\begin{equation*}
    \var{L} = \sum_{h=1}^n Q_h \var{q}_h.
\end{equation*}

\begin{example}[Caso in tre dimensioni]
    Consideriamo un punto materiale soggetto alla forza $\vF$ e calcoliamo il lavoro virtuale.
    \begin{gather*}
        \var{\vP} = \var{x}\vu{i} + \var{y}\vu{j} + \var{z}\vu{k}, \\
        \vF = F_x \vu{i} + F_y \vu{j} + F_z \vu{k}, \\
        \var{L} = \vF \cdot \var{\vP} = F_x \var{x} + F_y \var{y} + F_z \var{z}.
    \end{gather*}

    In questo caso le componenti lagrangiane della forza corrispondo a quelle lagrangiane. Inoltre, si tratta semplicemente di una forma differenziale lineare in $\var{x}, \var{y}, \var{z}$.
\end{example}

Le $Q_h$ sono funzioni
\begin{equation*}
    Q_h = Q_h(q_1, q_2, \dots, q_n; \dot{q}_1, \dot{q}_2, \dots, \dot{q}_n; t).
\end{equation*}

\begin{definition}[Sistema di forze posizionali]
    Un sistema di forze si dice \emph{posizionale} se le componenti lagrangiane dipendono solo dalle coordinate lagrangiane
    \begin{equation*}
        Q_h = Q_h(q_1, q_2, \dots, q_n).
    \end{equation*}
\end{definition}

In meccanica newotniana, se in un sistea di forze
\begin{equation*}
    \vF_1, \dots, \vF_N,
\end{equation*}
per ogni forza $\vF_i$ esiste una funzione scalare
\begin{equation*}
    V = V(x_1, y_1, z_1; x_2, y_2, z_2; \dots; x_N, y_N, z_N),
\end{equation*}
tale che
\begin{equation*}
    \vF_i = -\grad_i V \quad i = 1, 2, \dots, N,
\end{equation*}
allora il sistema di forze è conservativo.

\begin{definition}[Sistema di forze conservativo]
    Diremo che un sistema di forze agenti su $\mathcal{S}$ (a vincoli fissi) è \emph{conservativo} se esiste una funzione
    \begin{equation*}
        V=V(q_1, q_2, \dots, q_n),
    \end{equation*}
    detta \emph{energia potenziale}, tale che
    \begin{equation*}
        \var{L} = -\var{V},
    \end{equation*}
    ovvero
    \begin{equation*}
        Q_h = -\pdv{V}{q_h} \quad h = 1, 2, \dots, n,
    \end{equation*}
    dove le $Q_h$ sono le componenti lagrangiane della sollecitazione considerata.
\end{definition}

In un sistema a vincoli fissi il lavoro virtuale coincide con quello reale, per cui se il sistema di forze è conservativo, allora
\begin{equation*}
    \dd{L} = -\dd{V}.
\end{equation*}
Inoltre si ha che la variazione dell'energia cinetica è uguale al lavoro totale del sistema di forze,
\begin{equation*}
    \dd{T} = \dd{L}.
\end{equation*}
Quindi, se il sistema di forze è conservativo, si ha che
\begin{equation*}
    \dd{T} = \dd{L} + \dd{V} = 0,
\end{equation*}
per cui l'energia meccanica totale è costante:
\begin{equation*}
    E = T + V = \text{costante}.
\end{equation*}

Ora, nel sistema oltre alle forze che agiscono sui punti, ci sono anche le forze vincolari, per cui
\begin{equation*}
    \dd{T} = \underbrace{\dd{L}^{(a)}}_{\text{conservative}} + \dd{L}^{(v)},
\end{equation*}
dove $\dd{L}^{(a)}$ è il lavoro reale delle forze attive, mentre $\dd{L}^{(v)}$ è il lavoro reale delle forze vincolari. Se i vincoli sono fissi (quindi il sistema è scleronomo), lisci e bilaterali, allora  si ha
\begin{equation*}
    \begin{cases}
        \var{L}^{(v)} \geq 0 \\
        -\var{L}^{(v)} \geq 0
    \end{cases}
    \implies
    \begin{cases}
        \vF\cdot \var{\vP} \geq 0 \\
        -\vF\cdot \var{\vP} \geq 0
    \end{cases}
    \implies
    \vF\cdot \var{\vP} = 0,
\end{equation*}
per cui
\begin{equation*}
    \dd{T} = \dd{L}^{(a)} = -\dd{V}.
\end{equation*}

\subsection{Energia cinetica di un sistema olonomo}
L'energia cinetica di un sistema di $N$ punti materiali è data da
\begin{equation*}
    T = \frac{1}{2} \sum_{i=1}^N m_i v_i^2.
\end{equation*}
Teniamo però conto che possiamo esprimere ciascuna velocita come
\begin{equation*}
    \vv_i = \dv{\vP_i}{t} = \sum_{h=1}^n \pdv{\vP_i}{q_h} \dot{q}_h + \pdv{\vP_i}{t},
\end{equation*}
per cui otteniamo
\begin{align*}
    T &= \frac{1}{2} \sum_{i=1}^N m_i \left( \sum_{h=1}^n \pdv{\vP_i}{q_h} \dot{q}_h + \pdv{\vP_i}{t} \right)^2 \\
    &= \frac{1}{2} \sum_{i=1}^N m_i \left( \sum_{h=1}^n \pdv{\vP_i}{q_h} \dot{q}_h + \pdv{\vP_i}{t} \right)\cdot \left( \sum_{k=1}^n \pdv{\vP_i}{q_k} \dot{q}_k + \pdv{\vP_i}{t} \right) \\
    &= \frac{1}{2} \sum_{i=1}^N m_i \left[ \sum_{h,k=1}^n \pdv{\vP_i}{q_h} \cdot \pdv{\vP_i}{q_k} \dot{q}_h \dot{q}_k + 2 \sum_{h=1}^n \pdv{\vP_i}{q_h} \cdot \pdv{\vP_i}{t} \dot{q}_h + \left(\pdv{\vP_i}{t}\right)^2 \right] \\
    &= \underbrace{\frac{1}{2} \sum_{h,k=1}^n \left( \sum_{i=1}^N m_i \pdv{\vP_i}{q_h} \cdot \pdv{\vP_i}{q_k} \right) \dot{q}_h \dot{q}_k}_{T^2} + \underbrace{\sum_{h=1}^n \left( \sum_{i=1}^N m_i \pdv{\vP_i}{q_h} \cdot \pdv{\vP_i}{t} \right) \dot{q}_h}_{T^1} + \underbrace{\frac{1}{2} \sum_{i=1}^N m_i \left(\pdv{\vP_i}{t}\right)^2}_{T^0}\\
    &= T^2 + T^1 + T^0.
\end{align*}
Questo è un polinomio di secondo grado in $\dot{q}_i$. Definiamo
\begin{equation*}
    a_{hk} \defeq \sum_{i=1}^N m_i \pdv{\vP_i}{q_h} \cdot \pdv{\vP_i}{q_k}, \qquad a_{h} \defeq \sum_{i=1}^N m_i \pdv{\vP_i}{q_h} \cdot \pdv{\vP_i}{t},
\end{equation*}
per cui otteniamo
\begin{equation*}
    T_2 = \frac{1}{2} \sum_{h,k=1}^n a_{hk}\, \dot{q}_h \dot{q}_k, \qquad T_1 = \sum_{h=1}^n a_h \,\dot{q}_h, \qquad T_0 = \frac{1}{2} \sum_{i=1}^N m_i \left(\pdv{\vP_i}{t}\right)^2.
\end{equation*}

\begin{remark}
    Se il sistema è a vincoli indipendenti dal tempo, allora l'energia cinetica si riduce alla sola parte quadratica
    \begin{equation*}
        T = T^2.
    \end{equation*}

    Infatti si ha
    \begin{equation*}
        \pdv{\vP_i}{t} = 0 \quad \forall i = 1, 2, \dots, N,
    \end{equation*}
\end{remark}

Si ricorda che gli elementi $a_{hk}$ dipendono in generale dalle coordinate lagrangiane, per cui non si tratta di un termine costante.

La matrice $A = (a_{hk})$ è detta \emph{matrice di massa}

Proprietà: $T^2$ è una forma quadratica definita positiva.
\begin{proof}
    \begin{align*}
        \sum_{h,k=1}^n a_{hk}\, \xi_h \xi_k &= \sum_{h,k=1}^n\left(\sum_{i=1}^N m_i \pdv{\vP_i}{q_h}\cdot \pdv{\vP_i}{q_k}\right)\xi_h \xi_k \\
        &= \sum_{i=1}^N m_i \left( \sum_{h=1}^n \pdv{\vP_i}{q_h} \xi_h \right)^2 \geq 0,
    \end{align*}
    pertanto è semidefinita positiva. Resta da dimostrare che  è uguale a zero solo se $\xi_h = 0$ per ogni $h = 1, 2, \dots, n$. Se $T^2 = 0$, allora
    \begin{equation*}
        \sum_{h=1}^n \pdv{\vP_i}{q_h} \xi_h = 0 \quad \forall i = 1, 2, \dots, N.
    \end{equation*}

    Notiamo che la somma delle derivate $\pdv{\vP_i}{q_h}$ rappresenta una combinazione lineare delle colonne della matrice Jacobiana delle coordinate
    \begin{equation*}
        \left(\begin{array}{cccc}
            \pdv{x_1}{q_1} & \pdv{x_1}{q_2} & \dots & \pdv{x_1}{q_n} \\[0.2cm]
            \pdv{y_1}{q_1} & \pdv{y_1}{q_2} & \dots & \pdv{y_1}{q_n} \\[0.2cm]
            \pdv{z_1}{q_1} & \pdv{z_1}{q_2} & \dots & \pdv{z_1}{q_n} \\[0.2cm]
            \vdots & \vdots & \ddots & \vdots \\[0.2cm]
            \pdv{x_N}{q_1} & \pdv{x_N}{q_2} & \dots & \pdv{x_N}{q_n} \\[0.2cm]
            \pdv{y_N}{q_1} & \pdv{y_N}{q_2} & \dots & \pdv{y_N}{q_n} \\[0.2cm]
            \pdv{z_N}{q_1} & \pdv{z_N}{q_2} & \dots & \pdv{z_N}{q_n}
        \end{array}\right)
    \end{equation*}
    che ha rango $n$. Per cui la somma è uguale a zero solo se $\xi_h = 0$ per ogni $h = 1, 2, \dots, n$ in quanto le $q_i$ sono indipendenti.
\end{proof}